% Lösungen Einheit 4 von 5 – Steigungswinkel und Schnittwinkel (zusammenbau v0.1)
\einheitenkopf{Einheit 4 von 5 · Steigungswinkel und Schnittwinkel}

\erg{20}{a) ein Winkel – über den Tangens in die Steigung übersetzen, dann $f'(x) = m$ lösen. \quad b) $\alpha = \mathrm{tan}^{-1}(2) \approx 63{,}4^\circ$ \quad c) $f'(x) = x^2 - 2x$, $f'(3) = 3$; $\alpha = \mathrm{tan}^{-1}(3) \approx 71{,}6^\circ$ \quad d) $f(1) = \frac{3}{e} \approx 1{,}10$, $m = f'(1) = -\frac{1}{e} \approx -0{,}37$; Achsenabschnitt aus dem Punkt (nicht $f(1)$): $t(x) = -\frac{1}{e}x + \frac{4}{e} \approx -0{,}37x + 1{,}47$; Schnittwinkel mit der $x$-Achse $\mathrm{tan}^{-1}(\frac{1}{e}) \approx 20{,}2^\circ$ \quad e) $f'(x) = -\frac{3}{4}x$, $f'(-2) = 1{,}5$; Steigungswinkel $\approx 56{,}3^\circ$ gegen die Waagerechte; gegen die senkrechte Kante $90^\circ - 56{,}3^\circ \approx 33{,}7^\circ$ \quad f) $\alpha \approx -63{,}4^\circ$ zwischen der Tangente im Startpunkt $(0 | 4)$ und der Waagerechten (unterhalb); $\beta \approx 116{,}6^\circ$ ist der Winkel, den Podest und Rutschbahn am Übergang einschließen (der Knick) – nicht der Winkel zur Senkrechten. \quad g) $f'(0) = 0{,}5$; Steigungswinkel $\mathrm{tan}^{-1}(0{,}5) \approx 26{,}6^\circ$ und $\mathrm{tan}^{-1}(2) \approx 63{,}4^\circ$; Schnittwinkel $\approx 36{,}9^\circ$ (nicht $\mathrm{tan}^{-1}(2 - 0{,}5)$) \quad h) Schnittstelle $x = 0$; $f'(0) = 3k$, $g'(0) = -3k$; senkrecht: $3k \cdot (-3k) = -1$ (nicht die Symmetrie), also $k^2 = \frac{1}{9}$, $k = \frac{1}{3}$ \quad i) $f''(x) = 6x - 6 = 0$ bei $x = 1$, $W(1 | 0)$; $f'(1) = -3$; Wendetangente $t(x) = -3x + 3$, $y$-Achsenabschnitt $t(0) = 3$; Steigungswinkel $\mathrm{tan}^{-1}(-3) \approx -71{,}6^\circ$ (fällt) \quad j) $f'(x) = x^2 \ge \mathrm{tan}(60^\circ) \approx 1{,}73$; Grenze $x = \sqrt{\mathrm{tan}(60^\circ)} \approx 1{,}32$, also $1{,}32 \le x \le 3$ \quad k) $t$ hat den Steigungswinkel $45^\circ$, die Winkelhalbierende $22{,}5^\circ$ (nicht die halbe Steigung); Steigung von $g_u$: $\frac{f(u)}{u} = e^{-0{,}1u^2} = \mathrm{tan}(22{,}5^\circ) \approx 0{,}414$; $u^2 \approx 8{,}81$, $u \approx 2{,}97$ oder $u \approx -2{,}97$}
\erg{21}{a) Beim Kippen liegt der Boden als Tangente am Profil; ist diese Tangente die Gerade durch den Randpunkt, setzt der Rand auf – weiter lässt sich das Brett nicht neigen: der Winkel ist der größte Kippwinkel.}
\erg{22}{a) Pia zieht die Anstiege statt der Winkel ab. Richtig: $\mathrm{tan}^{-1}(3) - \mathrm{tan}^{-1}(1) \approx 71{,}6^\circ - 45^\circ \approx 26{,}6^\circ$. \quad b) Im Steigungsdreieck ist $m$ der Höhenunterschied geteilt durch den waagerechten Unterschied – Gegenkathete durch Ankathete des Winkels $\alpha$ an der Waagerechten, also $\mathrm{tan}\,\alpha$.}
